\documentclass[
    12pt,
    openright,
    oneside,
    a4paper,
    english,
    brazil
]{abntex2}

% ---
% PACOTES ESSENCIAIS
% ---
\usepackage[utf8]{inputenc}
\usepackage[T1]{fontenc}
\usepackage{lmodern}
\usepackage{indentfirst}
\usepackage{color}
\usepackage{xcolor}
\usepackage{graphicx}
\usepackage{microtype}
\usepackage{booktabs}
\usepackage{amsmath}
\usepackage{listings}
\usepackage{url}

% ---
% PACOTES DE CITAÇÃO E REFERÊNCIAS (ABNT)
% ---
\usepackage[alf,abnt-emphasize=bf,abnt-repeated-title-omission=yes]{abntex2cite}

% ---
% CONFIGURAÇÕES DO HYPERREF
% ---
\hypersetup{
    colorlinks=true,
    linkcolor=blue,
    citecolor=blue,
    filecolor=magenta,
    urlcolor=blue,
    bookmarksdepth=4
}

% ---
% CONFIGURAÇÃO DE CÓDIGO FONTE (LISTINGS)
% ---
\lstset{
    basicstyle=\ttfamily\footnotesize,
    keywordstyle=\color{blue}\bfseries,
    commentstyle=\color{gray}\itshape,
    stringstyle=\color{teal},
    numbers=left,
    numberstyle=\tiny\color{gray},
    stepnumber=1,
    numbersep=8pt,
    backgroundcolor=\color{white},
    showspaces=false,
    showstringspaces=false,
    showtabs=false,
    frame=single,
    tabsize=2,
    captionpos=b,
    breaklines=true,
    breakatwhitespace=false
}

% ---
% DADOS INSTITUCIONAIS (UTFPR CAMPUS GUARAPUAVA)
% ---
\instituicao{UNIVERSIDADE TECNOLÓGICA FEDERAL DO PARANÁ\par CÂMPUS GUARAPUAVA\par DEPARTAMENTO ACADÊMICO DE INFORMÁTICA\par CURSO SUPERIOR DE TECNOLOGIA EM SISTEMAS PARA INTERNET}

\titulo{PLATAFORMA EDUCACIONAL BASEADA EM SIMULAÇÃO POSIX NO NAVEGADOR PARA O ENSINO E AVALIAÇÃO PRÁTICA DE SISTEMAS OPERACIONAIS}

\autor{RICARDO MARTINS DE OLIVEIRA}

\local{Guarapuava}
\data{2026}

\orientador{Profª. Dra. Sediane Carmem Lunardi Hernandes}

\preambulo{Trabalho de Conclusão de Curso apresentado ao Curso Superior de Tecnologia em Sistemas para Internet do Câmpus Guarapuava da Universidade Tecnológica Federal do Paraná (UTFPR), como requisito parcial para a obtenção do título de Tecnólogo em Sistemas para Internet.}

% ---
% INÍCIO DO DOCUMENTO
% ---
\begin{document}
\selectlanguage{brazil}
\frenchspacing

% ----------------------------------------------------------
% ELEMENTOS PRÉ-TEXTUAIS
% ----------------------------------------------------------
\imprimircapa
\imprimirfolhaderosto*

% Resumo em Português
\setlength{\absparsep}{18pt}
\begin{resumo}
O aprendizado prático de Sistemas Operacionais (SO), em particular no ambiente Linux, constitui um pilar fundamental na formação de profissionais em Tecnologia da Informação. Tradicionalmente, as aulas práticas dependem do provisionamento de Máquinas Virtuais (VMs) locais ou de configurações em \textit{dual boot}. Entretanto, esse modelo impõe barreiras severas: elevado consumo de memória RAM e CPU, frequentes incompatibilidades de virtualização por \textit{hardware} (como desativação de VT-x na BIOS dos computadores pessoais dos alunos), tempo letivo substancial consumido pelo docente com suporte técnico de infraestrutura e reduzido engajamento discente frente a roteiros estáticos. Além disso, plataformas de simulação puramente executadas no cliente (\textit{client-side}) revelam-se inadequadas para avaliações somativas formais, uma vez que critérios de correção podem ser facilmente inspecionados ou adulterados via ferramentas de desenvolvedor (\textit{DevTools}). Este Trabalho de Conclusão de Curso apresenta a concepção e o desenvolvimento de uma plataforma educacional integrada para o ensino e avaliação prática de Sistemas Operacionais. A solução combina um núcleo de simulação Linux/POSIX de alta fidelidade desenvolvido em TypeScript puro com Sistema de Arquivos Virtual (\textit{Virtual File System} - VFS) executado de forma interativa no navegador, aliado a uma arquitetura de \textit{back-end} segura e desacoplada. O servidor é responsável pela gestão de turmas, controle temporal de exames mediante \textit{tokens} dinâmicos de liberação para aplicação presencial, agendamento de repescagens e segundas chamadas com autorizações restritas, e um motor de correção remota baseado em \textit{snapshots} assinados do estado do VFS. A validação técnica inicial conta com uma suíte de mais de 750 testes automatizados que cobrem chamadas POSIX e cenários reais de certificação. Conclui-se que a abordagem elimina a fricção de configuração local, otimiza o tempo didático em sala de aula e assegura a integridade pedagógica e avaliativa na disciplina de Sistemas Operacionais do curso de Tecnologia em Sistemas para Internet da UTFPR Câmpus Guarapuava.

\textbf{Palavras-chave:} Sistemas Operacionais. Linux. Simulação Web. Virtual File System. Educação em Computação. Avaliação Prática. UTFPR.
\end{resumo}


% Abstract em Inglês
\setlength{\absparsep}{18pt}
\begin{resumo}[Abstract]
\begin{otherlanguage*}{english}
Practical learning in Operating Systems (OS), particularly within the Linux ecosystem, is a cornerstone in the education of Information Technology professionals. Traditionally, hands-on lab sessions rely on local Virtual Machines (VMs) or dual-boot configurations. However, this model introduces severe barriers: high consumption of RAM and CPU resources, hardware virtualization incompatibilities (such as disabled VT-x in students' BIOS settings), significant instructional time lost to technical troubleshooting instead of core conceptual teaching, and diminished student engagement with static exercise guides. Furthermore, browser-based simulators running entirely on the client side prove unsuitable for formal summative assessments, as evaluation logic can be trivially inspected or altered via browser developer tools (DevTools). This Undergraduate Thesis presents the design and development of an integrated educational platform for practical learning and assessment in Operating Systems. The proposed solution combines a high-fidelity Linux/POSIX simulation engine built in pure TypeScript with an in-memory Virtual File System (VFS) running client-side with zero latency, paired with a secure, decoupled back-end architecture. The back-end orchestrates academic classrooms, time-controlled examinations enabled by dynamic release tokens during class hours, dedicated scheduling for make-up exams, and a server-side evaluation engine powered by signed VFS state snapshots. The initial technical validation comprises an automated test suite with over 750 passing tests covering POSIX compliance and industry certification scenarios. We conclude that this architecture removes local environment friction, maximizes active instructional time in the classroom, and guarantees assessment integrity for the Operating Systems curriculum at UTFPR Guarapuava Campus.

\textbf{Keywords:} Operating Systems. Linux. Web Simulation. Virtual File System. Computer Science Education. Practical Assessment. UTFPR.
\end{otherlanguage*}
\end{resumo}


% Lista de Abreviaturas e Siglas
\begin{siglas}
  \item[ABNT] Associação Brasileira de Normas Técnicas
  \item[ACM] \textit{Association for Computing Machinery}
  \item[API] \textit{Application Programming Interface} (Interface de Programação de Aplicação)
  \item[APT] \textit{Advanced Package Tool}
  \item[BIOS] \textit{Basic Input/Output System}
  \item[CI/CD] \textit{Continuous Integration / Continuous Delivery} (Integração e Entrega Contínuas)
  \item[CLI] \textit{Command-Line Interface} (Interface de Linha de Comando)
  \item[CPU] \textit{Central Processing Unit} (Unidade Central de Processamento)
  \item[DAINF] Departamento Acadêmico de Informática
  \item[DevTools] \textit{Developer Tools} (Ferramentas de Desenvolvedor no Navegador)
  \item[FHS] \textit{Filesystem Hierarchy Standard}
  \item[GID] \textit{Group Identifier} (Identificador de Grupo)
  \item[GNU] \textit{GNU's Not Unix}
  \item[GUI] \textit{Graphical User Interface} (Interface Gráfica do Usuário)
  \item[HTTP] \textit{Hypertext Transfer Protocol}
  \item[HTTPS] \textit{Hypertext Transfer Protocol Secure}
  \item[I/O] \textit{Input/Output} (Entrada e Saída)
  \item[IEEE] \textit{Institute of Electrical and Electronics Engineers}
  \item[JSON] \textit{JavaScript Object Notation}
  \item[JWT] \textit{JSON Web Token}
  \item[LLM] \textit{Large Language Model} (Modelo de Linguagem de Grande Escala)
  \item[LPIC] \textit{Linux Professional Institute Certification}
  \item[LTS] \textit{Long Term Support}
  \item[POSIX] \textit{Portable Operating System Interface}
  \item[RAM] \textit{Random Access Memory} (Memória de Acesso Aleatório)
  \item[REST] \textit{Representational State Transfer}
  \item[SBC] Sociedade Brasileira de Computação
  \item[SIGCSE] \textit{Special Interest Group on Computer Science Education}
  \item[SO] Sistemas Operacionais
  \item[SSH] \textit{Secure Shell}
  \item[TCC] Trabalho de Conclusão de Curso
  \item[TSI] Tecnologia em Sistemas para Internet
  \item[UID] \textit{User Identifier} (Identificador de Usuário)
  \item[UI] \textit{User Interface} (Interface com o Usuário)
  \item[UTFPR] Universidade Tecnológica Federal do Paraná
  \item[VFS] \textit{Virtual File System} (Sistema de Arquivos Virtual)
  \item[VM] \textit{Virtual Machine} (Máquina Virtual)
  \item[VT-x] \textit{Virtualization Technology} (Tecnologia de Virtualização Intel/AMD)
  \item[WSL] \textit{Windows Subsystem for Linux} (Subsistema Windows para Linux)
  \item[WSS] \textit{WebSocket Secure}
\end{siglas}


% Sumário
\pdfbookmark[0]{\contentsname}{toc}
\tableofcontents*
\cleardoublepage

% ----------------------------------------------------------
% ELEMENTOS TEXTUAIS
% ----------------------------------------------------------
\textual

\chapter{Introdução}
\label{cap:introducao}

A disciplina de Sistemas Operacionais (SO) integra o núcleo estruturante da formação científica e tecnológica em cursos superiores de Computação e Tecnologia da Informação. No âmbito do Curso Superior de Tecnologia em Sistemas para Internet (TSI) da Universidade Tecnológica Federal do Paraná (UTFPR) --- Câmpus Guarapuava, a compreensão prática sobre o funcionamento de núcleos de sistemas, gerenciamento de processos, sistemas de arquivos hierárquicos e modelos de permissões de acesso constitui uma competência determinante para a atuação dos discentes em desenvolvimento de \textit{software}, arquitetura de infraestrutura em nuvem e segurança da informação \cite{silberschatz2018operating, nemeth2017unix}.

Historicamente, o padrão POSIX (\textit{Portable Operating System Interface}) e as distribuições do ecossistema GNU/Linux estabeleceram-se como a espinha dorsal dos servidores web, contêineres de microsserviços e sistemas embarcados modernos \cite{posix2018, kerrisk2010linux}. Consequentemente, o domínio prático da Interface de Linha de Comando (\textit{Command-Line Interface} -- CLI) do Linux deixou de ser um conhecimento secundário para se tornar um requisito indispensável no perfil do egresso de TSI.

Entretanto, transpor esse conjunto teórico para a prática laboratorial efetiva impõe desafios logísticos, pedagógicos e tecnológicos que frequentemente comprometem a dinâmica pedagógica do docente e a curva de aprendizagem dos estudantes.

% ----------------------------------------------------------
\section{Contextualização e Dores Pedagógicas}
\label{sec:dores-pedagogicas}
% ----------------------------------------------------------

No contexto das atividades acadêmicas orientadas pela Profª. Dra. Sediane Carmem Lunardi Hernandes na UTFPR Câmpus Guarapuava, a condução das aulas práticas e das avaliações formais de Sistemas Operacionais esbarra em limitações materiais e operacionais recorrentes. Tais limitações podem ser categorizadas em quatro dimensões fundamentais:

\subsection{A Barreira das Máquinas Virtuais e a Heterogeneidade de Hardware}
A abordagem pedagógica convencional para práticas de laboratório baseia-se na instalação de Máquinas Virtuais (VMs), como VirtualBox ou VMware, ou na ativação do Subsistema Windows para Linux (WSL) nos computadores pessoais dos alunos. Essa estratégia exige computadores com capacidade substancial de memória volátil (RAM) e processamento, demandando tipicamente a alocação de no mínimo 2~GB a 4~GB de RAM para cada instância virtualizada, além de dezenas de \textit{gigabytes} de armazenamento em disco.

Na realidade dos discentes de graduação, uma fração expressiva possui computadores com limitações de \textit{hardware}. Além do consumo elevado de recursos, surgem barreiras de configuração em nível de \textit{firmware}: computadores com tecnologia de virtualização (Intel VT-x ou AMD-V) desativada de fábrica na BIOS/UEFI, conflitos de virtualização aninhada com plataformas antivírus e problemas de placas de rede em modo \textit{bridge} ou NAT.

\subsection{Tempo Letivo Desperdiçado com Suporte Técnico de Infraestrutura}
Como consequência direta da heterogeneidade e instabilidade dos ambientes locais, constata-se com frequência a perda de 30 a 45 minutos do início de cada aula prática exclusivamente no atendimento a chamados de suporte técnico elementar (``meu VirtualBox não inicializa'', ``minha máquina virtual perdeu a conexão com a rede'', ``meu teclado está desconfigurado no terminal''). 

Esse desperdício de tempo letivo subtrai a energia pedagógica do docente e a atenção da turma, desviando o foco do aprendizado de conceitos estruturais de Sistemas Operacionais --- como permissões de arquivos, mascaramento \textit{umask}, controle de usuários e manipulação de fluxos com \textit{pipes} e redirecionadores --- para a resolução de contingências de suporte de TI.

\subsection{O Desafio do Engajamento e a Fricção na Prática Autônoma}
A literatura em educação em computação evidencia que o aprendizado ativo e o construcionismo são potencializados quando os aprendizes experimentam ciclos rápidos de tentativa, erro e \textit{feedback} imediato \cite{papert1980mindstorms, benari2001constructivism}. Listas de exercícios estáticas em PDF ou apostilas teóricas sem validação imediata provocam desengajamento discente, pois o aluno, ao executar um comando incorreto no terminal, muitas vezes não compreende a causa da falha ou desiste por falta de orientação contextualizada.

\subsection{Vulnerabilidades e Fraudes em Avaliações Práticas Computadorizadas}
A aplicação de provas práticas individuais no ambiente de laboratório demanda garantias rigorosas de integridade avaliativa. Soluções de simuladores que operam exclusivamente no navegador do aluno (\textit{client-side}), sem uma arquitetura de retaguarda (\textit{back-end}) segura, apresentam vulnerabilidades críticas:
\begin{enumerate}
    \item \textbf{Inspeção e Adulteração via DevTools:} Qualquer estudante dotado de conhecimentos em desenvolvimento web pode abrir as ferramentas de depuração do navegador (F12 / \textit{Developer Tools}), inspecionar as rotinas de validação em JavaScript, descobrir as condições exatas esperadas pelos testes e alterar as variáveis de estado em memória para obter nota máxima sem resolver o desafio;
    \item \textbf{Controle Temporal e Presença em Sala:} A ausência de um mecanismo de liberação imediata por \textit{tokens} ou senhas dinâmicas impede o professor de sincronizar o início da avaliação presencial, abrindo brechas para acessos antecipados ou compartilhamento indevido de provas;
    \item \textbf{Gestão de Repescagens e Segundas Chamadas:} Em um semestre letivo universitário, é comum que alunos necessitem de avaliações em regime de segunda chamada ou repescagem por motivos de saúde ou faltas justificadas. A plataforma precisa possibilitar a criação de janelas temporais personalizadas com credenciais restritas a alunos específicos, sem expor a avaliação à turma regular.
\end{enumerate}

% ----------------------------------------------------------
\section{Justificativa}
\label{sec:justificativa}
% ----------------------------------------------------------

Diante do diagnóstico apresentado, justifica-se o desenvolvimento de uma plataforma educacional que elimine a dependência de ambientes de virtualização pesados e, simultaneamente, supra a lacuna de integridade nas avaliações práticas.

A proposta de transição do protótipo funcional de simulador em TypeScript puro para uma plataforma educacional completa permite unir o melhor de dois mundos:
\begin{itemize}
    \item \textbf{No cliente (\textit{front-end}):} Um simulador de alta fidelidade ao padrão POSIX com Sistema de Arquivos Virtual (\textit{Virtual File System} -- VFS), que roda com latência zero diretamente no navegador web, dispensando instalações locais, configurações de BIOS e alocações massivas de memória;
    \item \textbf{No servidor (\textit{back-end}):} Uma arquitetura desacoplada capaz de gerenciar turmas e alunos, gerar \textit{tokens} temporais de prova, controlar janelas de aplicação para turmas regulares e repescagens, e executar o motor de validação do estado do VFS de forma isolada, impedindo fraudes originadas no cliente e fornecendo relatórios diagnósticos de aprendizagem para a docente.
\end{itemize}

Dessa forma, a pesquisa e a implementação técnica resultam em um produto tecnológico de impacto direto na qualidade do ensino de Sistemas Operacionais na UTFPR, com potencial de extensão a outras instituições de ensino público e privado.

% ----------------------------------------------------------
\section{Objetivos}
\label{sec:objetivos}
% ----------------------------------------------------------

\subsection{Objetivo Geral}
Desenvolver e avaliar uma plataforma educacional integrada para o ensino e a avaliação prática de Sistemas Operacionais, composta por um simulador Linux/POSIX interativo executado no navegador e um serviço de \textit{back-end} seguro para gestão acadêmica, controle temporal de exames e motor de correção de VFS desacoplado contra fraudes.

\subsection{Objetivos Específicos}
Para o cumprimento do objetivo geral, estabelecem-se os seguintes objetivos específicos:
\begin{enumerate}
    \item Mapear e formalizar os contratos de API e o modelo de dados relacional para representação de usuários, turmas, avaliações, janelas temporais de aplicação e submissões com telemetria;
    \item Projetar um mecanismo de serialização e asserção de estado de Sistema de Arquivos Virtual (VFS), permitindo que o \textit{back-end} receba \textit{snapshots} do cliente e execute a correção de tarefas laboratoriais em ambiente protegido;
    \item Implementar o protocolo de segurança para aplicação de avaliações em tempo real mediante \textit{tokens} dinâmicos de liberação para turmas presenciais e controle de janelas específicas para repescagens e segundas chamadas;
    \item Integrar a interface do simulador existente (VFS, parser Bash e emulação de terminal) aos serviços de autenticação, carregamento de exercícios guiados e submissão de exames da nova plataforma;
    \item Validar a fidelidade da simulação e a integridade da plataforma através de testes automatizados unitários e de integração, além de analisar a aplicabilidade da ferramenta no contexto pedagógico da disciplina de Sistemas Operacionais da UTFPR Câmpus Guarapuava.
\end{enumerate}

% ----------------------------------------------------------
\section{Estrutura do Trabalho}
\label{sec:estrutura}
% ----------------------------------------------------------

Este documento está estruturado em seis capítulos:
\begin{itemize}
    \item O \textbf{Capítulo \ref{cap:introducao}} contextualiza a pesquisa, elenca as dores pedagógicas vivenciadas no ensino prático de Sistemas Operacionais, estabelece a justificativa do projeto e delineia os objetivos geral e específicos;
    \item O \textbf{Capítulo 2 (Fundamentação Teórica)} aborda os princípios teóricos de arquitetura de sistemas operacionais, o padrão POSIX, a hierarquia FHS, a semântica de chamadas de sistema, bem como os referenciais pedagógicos de ensino de computação e aprendizagem ativa;
    \item O \textbf{Capítulo 3 (Arquitetura da Plataforma)} detalha o projeto arquitetural da solução, os modelos de dados, o fluxo de comunicação entre cliente e servidor, a especificação das APIs REST e a estratégia de validação desacoplada de VFS;
    \item O \textbf{Capítulo 4 (Implementação)} descreve os aspectos técnicos da engenharia do \textit{software}, englobando o motor de simulação in-memory, o parser de comandos, a construção dos serviços de retaguarda e as camadas de segurança de prova;
    \item O \textbf{Capítulo 5 (Avaliação e Resultados)} apresenta os testes automatizados, a verificação da fidelidade aos comandos Linux, a resistência contra tentativas de manipulação em cliente e a análise comparativa frente a soluções baseadas em VMs;
    \item O \textbf{Capítulo 6 (Conclusão e Trabalhos Futuros)} sintetiza as contribuições alcançadas pelo trabalho acadêmico, discute as limitações remanescentes e sugere direções para extensões futuras da plataforma.
\end{itemize}

\chapter{Fundamentação Teórica}
\label{cap:fundamentacao}

Este capítulo apresenta a base conceitual da plataforma. A primeira parte trata do lado técnico: a arquitetura de sistemas operacionais, o padrão POSIX e a organização do sistema de arquivos no Linux. A segunda parte trata do lado pedagógico: a Teoria da Carga Cognitiva, a Aprendizagem Ativa, o \textit{Microlearning} e a dispersão causada pelos \textit{smartphones} em sala de aula.

% ----------------------------------------------------------
\section{Arquitetura de Sistemas Operacionais e o Padrão POSIX}
\label{sec:posix}
% ----------------------------------------------------------

Um sistema operacional é a camada que fica entre o \textit{hardware} e os programas do usuário \cite{silberschatz2018operating, tanenbaum2015modern}. Segundo \citeonline{tanenbaum2015modern}, ele cumpre dois papéis. O primeiro é o de \textbf{gerenciador de recursos}: distribui processador, memória, discos e dispositivos de entrada e saída entre programas que competem por eles. O segundo é o de \textbf{máquina estendida} (\textit{extended machine}): oferece abstrações convenientes, como arquivos, processos e \textit{sockets}, que escondem os detalhes do \textit{hardware}.

Os programas em espaço de usuário (\textit{user space}) conversam com o núcleo (\textit{kernel space}) por meio das chamadas de sistema (\textit{system calls} ou \textit{syscalls}). Com a multiplicação de variantes proprietárias do Unix, essas interfaces divergiram e a compatibilidade entre sistemas se perdeu. Para resolver o problema, o IEEE, a ISO e \textit{The Open Group} criaram o padrão POSIX (\textit{Portable Operating System Interface}) \cite{posix2018}.

O POSIX define contratos uniformes para operações de Entrada e Saída baseadas em descritores de arquivo, para criação e controle de processos (\texttt{fork}, \texttt{exec}, \texttt{wait}), para sinais assíncronos e para o comportamento do interpretador de comandos (\textit{shell}) e dos utilitários de linha de comando \cite{kerrisk2010linux}. Quem domina essas interfaces consegue escrever programas portáveis entre as diferentes distribuições Linux e os demais sistemas Unix-like.

% ----------------------------------------------------------
\section{Sistemas de Arquivos e o Padrão FHS}
\label{sec:fhs}
% ----------------------------------------------------------

No Unix e no Linux vale o princípio de que ``tudo é um arquivo'' (\textit{everything is a file}). Não só os dados em disco, mas também dispositivos, periféricos, canais de comunicação entre processos e estruturas do próprio núcleo aparecem como nós de uma única árvore de diretórios, com raiz em \texttt{/} \cite{nemeth2017unix, kerrisk2010linux}.

Para que cada distribuição não organize essa árvore à sua maneira, a \textit{Linux Foundation} mantém o \textit{Filesystem Hierarchy Standard} (FHS) \cite{fhs30}. O Quadro~\ref{qua:fhs} resume os principais diretórios definidos pelo padrão.

\begin{quadro}[htbp]
    \caption{Principais diretórios definidos pelo FHS}
    \label{qua:fhs}
    \footnotesize
    \begin{tabularx}{\textwidth}{|p{3.4cm}|L|p{3.6cm}|}
        \hline
        \textbf{Diretório} & \textbf{Finalidade} & \textbf{Exemplo} \\
        \hline
        \texttt{/bin}, \texttt{/sbin} & Binários essenciais para operar e administrar o sistema & \texttt{ls}, \texttt{mount} \\
        \hline
        \texttt{/etc} & Configurações globais da máquina e dos serviços & \texttt{/etc/passwd} \\
        \hline
        \texttt{/home} & Dados pessoais dos usuários comuns & \texttt{/home/ana} \\
        \hline
        \texttt{/root} & Diretório pessoal do superusuário & --- \\
        \hline
        \texttt{/var} & Dados que mudam durante a operação: \textit{logs}, bancos de dados, filas & \texttt{/var/log/syslog} \\
        \hline
        \texttt{/opt}, \texttt{/usr/local} & \textit{Software} adicional e programas compilados localmente & \texttt{/usr/local/bin} \\
        \hline
        \texttt{/tmp} & Arquivos temporários, protegidos pelo \textit{sticky bit} & \texttt{/tmp/sessao.lock} \\
        \hline
        \texttt{/dev}, \texttt{/proc}, \texttt{/sys} & Sistemas de arquivos virtuais que expõem dispositivos e estruturas do \textit{kernel} & \texttt{/dev/sda}, \texttt{/proc/cpuinfo} \\
        \hline
    \end{tabularx}
    \fonte{Adaptado de \citeonline{fhs30}.}
\end{quadro}

A segurança do sistema de arquivos se apoia em identidades de usuário (\textit{User ID} -- UID), em grupos primários e secundários (\textit{Group ID} -- GID) e na matriz de permissões POSIX: leitura (\textit{r} = 4), escrita (\textit{w} = 2) e execução (\textit{x} = 1), aplicadas ao dono (\textit{user}), ao grupo (\textit{group}) e aos demais (\textit{others}) \cite{nemeth2017unix}. Há ainda atributos especiais como \textit{SetUID}, \textit{SetGID} e o \textit{sticky bit} (octal 1000). Este último é o que impede, em diretórios compartilhados como \texttt{/tmp}, que um usuário apague arquivos criados por outro.

% ----------------------------------------------------------
\section{A Teoria da Carga Cognitiva no Ensino de Computação}
\label{sec:carga-cognitiva}
% ----------------------------------------------------------

A Teoria da Carga Cognitiva (\textit{Cognitive Load Theory} -- CLT), proposta por \citeonline{sweller1988cognitive} e ampliada nas décadas seguintes \cite{sweller2011cognitive, paas2003cognitive}, parte da distinção entre memória de trabalho (\textit{working memory}) e memória de longo prazo (\textit{long-term memory}).

A memória de longo prazo tem capacidade praticamente ilimitada para guardar estruturas de conhecimento, os chamados esquemas mentais. A memória de trabalho, ao contrário, é um gargalo: consegue manter e manipular poucos elementos novos ao mesmo tempo, e por pouco tempo \cite{sweller1988cognitive}. Quando uma atividade exige mais do que ela suporta, ocorre sobrecarga cognitiva (\textit{cognitive overload}) e o aprendizado não acontece.

A literatura da área \cite{paas2003cognitive, sweller2011cognitive} divide o esforço mental em três tipos de carga, descritos no Quadro~\ref{qua:cargas-cognitivas} com exemplos do ensino de Sistemas Operacionais.

\begin{quadro}[htbp]
    \caption{Os três tipos de carga cognitiva aplicados ao ensino de SO}
    \label{qua:cargas-cognitivas}
    \footnotesize
    \begin{tabularx}{\textwidth}{|p{2.4cm}|L|L|}
        \hline
        \textbf{Carga} & \textbf{O que é} & \textbf{Exemplo no ensino de SO} \\
        \hline
        \textbf{Intrínseca} & Complexidade do próprio conteúdo e do número de elementos que interagem entre si & Entender a precedência das permissões octais, a sincronização de processos ou o encadeamento de \textit{pipes} \\
        \hline
        \textbf{Germana} & Esforço produtivo de construir e consolidar esquemas na memória de longo prazo & Relacionar \texttt{chmod}, \texttt{umask} e \texttt{ls -l} em um modelo mental único de permissões \\
        \hline
        \textbf{Extrínseca} & Esforço desperdiçado com a forma de apresentação ou com obstáculos que não ensinam nada & Configurar o hipervisor, ativar VT-x na BIOS, depurar a rede da máquina virtual \\
        \hline
    \end{tabularx}
    \fonte{Autoria própria (2026), com base em \citeonline{paas2003cognitive} e \citeonline{sweller2011cognitive}.}
\end{quadro}

No ensino prático tradicional, pedir a iniciantes que instalem hipervisores, ajustem a BIOS/UEFI e consertem a rede da VM gera uma grande \textbf{carga extrínseca} \cite{silva2021ensino}. Essa carga ocupa a memória de trabalho e esgota a atenção do aluno antes que ele chegue aos conceitos de sistemas operacionais propriamente ditos.

Por isso, uma ferramenta educacional deve começar eliminando a carga extrínseca, oferecendo um ambiente que esteja pronto para uso imediato \cite{guo2013online}.

% ----------------------------------------------------------
\section{Aprendizagem Ativa, Microlearning e Ambientes Web Interativos}
\label{sec:aprendizagem-ativa}
% ----------------------------------------------------------

O ensino expositivo, baseado na transmissão passiva de conteúdo e em longas apostilas em PDF, tem eficácia limitada para desenvolver habilidades práticas em computação \cite{benari2001constructivism, freeman2014active}.

Em uma meta-análise publicada nos \textit{Proceedings of the National Academy of Sciences}, \citeonline{freeman2014active} compararam centenas de estudos e mostraram que a Aprendizagem Ativa (\textit{Active Learning}) em disciplinas de ciências, tecnologia, engenharia e matemática (STEM) eleva as notas médias e reduz em 55\% as taxas de reprovação em relação às aulas puramente expositivas.

O \textit{Microlearning}, por sua vez, propõe dividir conteúdos complexos em pequenas unidades (\textit{micro-lessons}), cada uma com um único objetivo prático e imediato \cite{kapp2019microlearning}. Cada lição combina uma explicação curta com um desafio que se resolve na hora, criando ciclos rápidos de tentativa, erro e acerto.

Ambientes interativos que rodam no próprio navegador, sem instalar nada, têm se destacado na literatura de Educação em Computação (SIGCSE/ACM) como forma de aumentar o engajamento e democratizar o acesso \cite{guo2013online, silva2021ensino}. Como mostra \citeonline{guo2013online}, remover a etapa de instalação de compiladores e interpretadores reduz drasticamente o atrito inicial e aumenta o tempo dedicado à experimentação.

% ----------------------------------------------------------
\section{O Perfil dos Novos Discentes e a Dispersão Atencional por Smartphones}
\label{sec:dispersao-smartphones}
% ----------------------------------------------------------

Quem ingressa no ensino superior nesta década pertence a uma geração cujos hábitos de atenção foram moldados pelos \textit{smartphones} e pela conexão permanente à internet \cite{kuznekoff2013impact}. Esse público consome informação de forma fragmentada, espera respostas imediatas dos sistemas e tende à multitarefa digital.

Estudos sobre o uso de celulares em sala de aula no ensino superior mostram que a distração digital prejudica a retenção de conteúdo e o engajamento, sobretudo quando a atividade perde dinamismo \cite{kuznekoff2013impact}. Em um laboratório, o atrito técnico excessivo, como esperar a máquina virtual iniciar ou encarar uma tela congelada, funciona como gatilho: o estudante deixa o computador de lado e pega o celular.

Depois que a atenção migra para redes sociais e mensagens, recuperá-la custa caro, e o aluno perde o ritmo da aula \cite{kuznekoff2013impact, silva2021ensino}. Ambientes com resposta visual instantânea e terminal interativo no navegador fazem o caminho inverso: prendem a atenção do estudante na atividade desde os primeiros minutos.

\chapter{Arquitetura da Plataforma}
\label{cap:arquitetura}

Este capítulo detalha a concepção arquitetural da plataforma educacional, demonstrando como a transição do protótipo inicial executado puramente no cliente para um modelo híbrido e reativo integrado cliente-servidor viabiliza avaliações seguras, supervisão pedagógica em tempo real, suporte síncrono em sala de aula e escalabilidade no ambiente acadêmico.

% ----------------------------------------------------------
\section{Visão Geral da Arquitetura Híbrida e Reativa}
\label{sec:visao-arquitetural}
% ----------------------------------------------------------

O projeto arquitetural da plataforma foi concebido para equacionar dois requisitos aparentemente antagônicos: por um lado, a necessidade de proporcionar uma experiência de terminal com \textbf{latência zero e alta responsividade} para o estudante; por outro lado, a exigência indispensável de \textbf{segurança avaliativa, integridade contra fraudes e supervisão docente contínua}.

No modelo convencional de aplicações web educacionais baseadas em requisições HTTP REST síncronas, a validação de tarefas práticas impõe atritos significativos ao usuário: o estudante precisa interromper seu fluxo de digitação no terminal, mudar o foco para a interface gráfica e acionar manualmente botões como ``Enviar Resposta'' ou ``Validar Tarefa''. Esse paradigma fragmenta a concentração do aprendiz e introduz latências de rede indesejadas a cada tentativa.

Para superar essa limitação, a plataforma adota uma **arquitetura reativa bidirecional orientada a eventos**, sustentada por canais persistentes seguros via \textit{WebSockets} (\textit{WSS}), operando em conformidade com o seguinte arranjo de responsabilidades:

\begin{itemize}
    \item \textbf{Camada de Execução Interativa (Cliente / Navegador):} O núcleo do simulador POSIX, a pilha de sessões Bash e a árvore do Sistema de Arquivos Virtual (\textit{VFS}) operam em memória local no navegador do aluno, desenvolvidos em TypeScript puro. Essa decisão assegura que a digitação de caracteres, a navegação pelo histórico, o autocompletar via tecla \texttt{Tab} e a renderização de cores ANSI ocorram instantaneamente, sem gerar tráfego de rede e sem sobrecarregar os servidores institucionais com a instanciação de contêineres individuais por aluno;
    \item \textbf{Camada de Validação e Orquestração (Servidor / Back-end):} O servidor atua como autoridade central confiável (\textit{Trusted Authority}). As regras de correção, os gabaritos das questões, os pesos das avaliações e as rotinas de asserção residem com exclusividade no servidor, tornando-os completamente inacessíveis à inspeção discente via console de desenvolvedor (\textit{DevTools});
    \item \textbf{Canal de Supervisão Docente (Cockpit do Professor):} O fluxo contínuo de eventos transmitidos via \textit{WebSockets} alimenta um painel administrativo em tempo real para o docente. A professora tem a capacidade de monitorar o progresso simultâneo de todos os terminais da turma durante uma prova prática presencial, identificar imediatamente discentes estagnados em mensagens de erro recorrentes e intervir pedagogicamente por meio de mensagens direcionadas ou avisos sincronizados.
\end{itemize}

% ----------------------------------------------------------
\section{Modelo de Dados Relacional}
\label{sec:modelo-dados}
% ----------------------------------------------------------

A gestão pedagógica e a persistência de longo prazo assentam-se sobre um banco de dados relacional estruturado (PostgreSQL), modelado para contemplar as entidades fundamentais do ecossistema acadêmico:
\begin{itemize}
    \item \textbf{Usuários e Perfis:} Diferenciação rigorosa entre os níveis de acesso Público, Aluno (vinculado por Registro Acadêmico - RA) e Professor;
    \item \textbf{Turmas e Inscrições:} Agrupamento semestral de turmas da disciplina de Sistemas Operacionais, possibilitando a matrícula e o acompanhamento de turmas regulares e especiais;
    \item \textbf{Avaliações e Janelas de Aplicação:} Modelagem de provas práticas associadas a janelas temporais estritas e chaves dinâmicas de liberação (\textit{tokens}), garantindo suporte nativo a repescagens e segundas chamadas isoladas;
    \item \textbf{Sessões e Submissões:} Registro auditável de cada tentativa individual de prova, armazenando o carimbo de data/hora validado pelo servidor, a nota obtida e a trilha completa de telemetria dos comandos executados.
\end{itemize}

% ----------------------------------------------------------
\section{Protocolo de Validação Desacoplada de VFS e Telemetria em Tempo Real}
\label{sec:validacao-desacoplada}
% ----------------------------------------------------------

Para mitigar integralmente os riscos de fraude via ferramentas de desenvolvedor (\textit{DevTools}) e, simultaneamente, eliminar a necessidade de botões manuais de envio que quebram o fluxo de raciocínio do discente, a validação de cada desafio opera através de um fluxo assíncrono de eventos via \textit{WebSockets}.

\subsection{Fluxo Reativo de Feedback Instantâneo}
A interação entre o terminal local e o motor de avaliação em retaguarda segue um ciclo não-bloqueante de três etapas:
\begin{enumerate}
    \item \textbf{Execução e Emissão em Background:} A cada quebra de linha (\texttt{Enter}) no terminal, o simulador executa a instrução instantaneamente no VFS local. Paralelamente à exibição imediata do resultado visual na tela do aluno, um emissor em segundo plano despacha pelo canal \textit{WebSocket} um \textit{payload} leve contendo o comando digitado, o código de retorno (\textit{exit code}), a identificação da questão e a árvore de nós ou metadados alterados (\textit{snapshot} diferencial);
    \item \textbf{Asserção em Ambiente Seguro:} O motor de avaliação no \textit{back-end} reconstitui o estado do VFS em memória protegida e aplica a bateria de testes de aceitação correspondente à tarefa ativa. Como as regras de teste residem unicamente no servidor, qualquer tentativa de manipulação de variáveis locais no navegador é inócua;
    \item \textbf{Retorno de Evento e Conclusão Automática:} Se as asserções forem plenamente satisfeitas, o servidor emite de volta pelo socket um evento assíncrono \texttt{task:completed}. A interface do aluno intercepta esse evento e exibe instantaneamente o \textit{feedback} visual de sucesso (animação de conclusão e avanço no roteiro), mantendo o aprendiz em estado contínuo de imersão e foco cognitivo.
\end{enumerate}

\subsection{Natureza Multifacetada do Linux e Auditoria Comportamental}
Essa abordagem fundamenta-se na própria natureza do ecossistema Linux, no qual inexiste um comando único para resolver determinada tarefa operacional. Por exemplo, a atribuição de permissões de leitura e escrita a um arquivo pode ser executada por notação octal (\texttt{chmod 644 arquivo.txt}) ou simbólica (\texttt{chmod u=rw,go=r arquivo.txt}); a busca por padrões de texto pode ser conduzida combinando \texttt{cat} e \texttt{grep} via \textit{pipes} ou diretamente através do argumento de entrada do utilitário \texttt{grep}.

Dessa forma, o motor de retaguarda atua em dois níveis complementares:
\begin{itemize}
    \item \textbf{Validação por Estado (Determinística):} O servidor avalia exclusivamente o resultado final refletido no VFS (verificando se o arquivo existe no caminho esperado, se possui o UID correto e se as permissões atendem à máscara exigida), assegurando aprovação pedagógica independentemente da rota sintática adotada pelo discente;
    \item \textbf{Análise Comportamental e Auditoria em Tempo Real:} A persistência cronológica de todas as instruções permite à docente mapear a proficiência dos comandos utilizados, diagnosticar lacunas conceituais recorrentes e garantir a integridade formal da avaliação contra tentativas de trapaça.
\end{itemize}

% ----------------------------------------------------------
\section{Governança Avaliativa, Telemetria de Segurança e Políticas Anti-Fraude}
\label{sec:governanca-seguranca}
% ----------------------------------------------------------

A aplicação de avaliações práticas em ambiente computacional exige um arcabouço rigoroso de governança que equilibre a supervisão pedagógica em tempo real, a prevenção ativa contra tentativas de trapaça e a tolerância a comportamentos acidentais. Em consonância com a arquitetura reativa orientada a eventos, a plataforma estrutura seus mecanismos de controle em cinco eixos operacionais: o suporte síncrono em sala de aula, a detecção heurística de ferramentas de inspeção, o bloqueio de concorrência de sessões, a telemetria forense de rede e a gestão de janelas temporais com chaves dinâmicas.

\subsection{Dinâmica Pedagógica da Fila de Suporte: A Mãozinha Virtual}
\label{sec:fila-suporte}

Em aulas laboratoriais práticas com turmas numerosas, um dos gargalos pedagógicos mais críticos reside na gestão do atendimento a dúvidas: múltiplos alunos erguem os braços simultaneamente, o docente perde a ordem cronológica dos chamados e dúvidas idênticas acabam sendo respondidas individualmente de forma exaustiva e repetitiva. Para solucionar esse problema, a plataforma incorpora o protocolo síncrono da \textbf{Fila de Suporte (Mãozinha Virtual)}, estruturado em três etapas operacionais integradas via \textit{WebSockets}:

\begin{enumerate}
    \item \textbf{Abertura Qualificada do Pedido de Ajuda (Levantar a Mão):} Ao acionar o comando de suporte na interface do terminal (``Pedir Ajuda''), o sistema exige obrigatoriamente que o estudante preencha um resumo textual sucinto descrevendo a natureza da sua dúvida ou o obstáculo enfrentado. Essa exigência pedagógica exercita a capacidade de síntese e metacognição do discente, forçando-o a estruturar o problema conceitualmente e eliminando chamados vazios ou acidentais. Uma vez despachado o evento \texttt{help:request}, o servidor insere um novo cartão na fila do painel docente (\textit{Cockpit}), ordenado pela disciplina cronológica FIFO (\textit{First-In, First-Out}), com data/hora, identificação do discente, terminal e descrição;
    \item \textbf{Cancelamento Autônomo (Abaixar a Mão):} Durante a espera na fila, caso o estudante consiga resolver o problema sozinho ou acompanhar uma explicação coletiva em sala, ele pode acionar a qualquer momento a opção ``Abaixar a Mão'', despachando \texttt{help:cancel}. O servidor remove instantaneamente o cartão da fila docente via \textit{WebSocket}, valorizando a autorregulação e evitando atendimentos desnecessários;
    \item \textbf{Atendimento Privado e Broadcast Inteligente:} A professora dispõe de dois modos de atendimento: a resposta individual discreta (diálogo 1:1 sobreposto ao terminal do aluno) e a resposta em lote por transmissão pública (\textit{multiselect}). Na resposta em lote, a docente seleciona múltiplos chamados correlatos e redige uma única orientação, a qual é transmitida para a turma mencionando os estudantes contemplados (exemplo: ``\textit{Respondendo às dúvidas de Aluno A, Aluno B e Aluno C: [Orientações da Docente]}''), resolvendo simultaneamente todos os chamados agrupados.
\end{enumerate}

\subsection{Políticas de Integridade e Detecção de DevTools com Avisos Escalonados}
\label{sec:integridade-devtools}

A integridade de uma avaliação prática computadorizada demanda proteção ativa contra a inspeção e a manipulação do código de execução no cliente. No entanto, em ambientes educacionais, o disparo de atalhos de teclado (como \texttt{F12}, \texttt{Ctrl+Shift+I} ou a alternância rápida de janelas com \texttt{Alt+Tab}) pode ocorrer de forma involuntária ou acidental pelo estudante. Para mitigar esse atrito sem prescindir do rigor avaliativo, o sistema implementa uma política de avisos escalonados e detecção em tempo real:

\begin{enumerate}
    \item \textbf{Sensoriamento Contínuo no Cliente:} O cliente mantém sentinelas em segundo plano que monitoram variações anômalas no redimensionamento da janela de exibição (\textit{window inner/outer dimensions mismatch}), limiares de tempo de execução causados por pontos de interrupção (\textit{debugger timing traps}) e eventos de perda de visibilidade da página por meio da API W3C \texttt{visibilitychange};
    \item \textbf{Primeira Ocorrência (Advertência Pedagógica):} Na primeira detecção de abertura de ferramentas de depuração ou perda de foco prolongada da aba de prova, o sistema não penaliza imediatamente a nota. Em vez disso, bloqueia temporariamente a digitação e apresenta um modal instrutivo em tela cheia, advertindo formalmente o discente sobre as regras de integridade e registrando a advertência local;
    \item \textbf{Reincidência e Notificação ao Painel Docente:} Caso o comportamento persista ou o estudante reabra as ferramentas de desenvolvedor, o cliente despacha imediatamente o evento \texttt{security:violation} pelo canal \textit{WebSocket}. O servidor registra a infração na trilha de telemetria forense e atualiza o \textit{card} do aluno no painel do professor com um indicador visual de alerta crítico (\textit{badge} de atenção). O sistema preserva a soberania da autoridade pedagógica, conferindo à docente a prerrogativa de verificar os registros circunstanciais, advertir presencialmente o aluno, desconsiderar falsos positivos ou anular a tentativa avaliativa.
\end{enumerate}

\subsection{Bloqueio de Concorrência de Sessão (Mutex Anti-Proxy)}
\label{sec:mutex-sessao}

Uma vulnerabilidade recorrente em exames computadorizados é o compartilhamento indevido de credenciais de acesso, no qual um discente presente no laboratório cede seu login a um colaborador externo (\textit{proxy} humano) para a resolução remota paralela da prova.

Para erradicar essa prática, a plataforma implementa uma trava de exclusão mútua de sessão (\textit{Session Mutex}) gerenciada no servidor sobre o protocolo de \textit{WebSockets}:
\begin{itemize}
    \item \textbf{Unicidade Estrita de Conexão:} Cada estudante autenticado por Registro Acadêmico (RA) só pode manter exatamente uma conexão \textit{socket} ativa por avaliação;
    \item \textbf{Derrubada Automática e Registro de Colisão:} Se uma nova tentativa de conexão for estabelecida com as mesmas credenciais (seja a partir de outro navegador, aba em modo anônimo ou dispositivo remoto distinto), o servidor derruba instantaneamente o canal de comunicação anterior com o código de encerramento \texttt{4409 Conflict} e emite a notificação: ``\textit{Conexão Encerrada: Nova sessão de avaliação detectada em outro dispositivo}'';
    \item \textbf{Alerta Imediato de Fraude:} O evento de colisão gera uma notificação prioritária no painel docente em tempo real, informando o endereço IP de ambas as pontas e o carimbo de data/hora da tentativa, permitindo a flagrância do compartilhamento indevido.
\end{itemize}

\subsection{Telemetria Forense de Rede, Geolocalização e Limitações de Sandbox}
\label{sec:telemetria-forense}

Como camada complementar de auditoria pós-exame, cada submissão e evento crítico de sessão é enriquecido com metadados de contexto de rede e ambiente operacional:
\begin{itemize}
    \item \textbf{Captura e Triagem de Endereços IP:} O servidor extrai o endereço IP público de origem e analisa o cabeçalho de proxy reverso corporativo. Em exames presenciais, o sistema permite à docente configurar um filtro de sub-rede institucional, confirmando automaticamente se todas as conexões provêm do bloco de rede e do \textit{gateway} exclusivo do laboratório físico da UTFPR;
    \item \textbf{Fingerprinting de Navegador e Coordenadas:} O cliente realiza uma composição de atributos de ambiente (\textit{User-Agent}, resolução de vídeo, idioma do sistema e \textit{timezone}), associando um \textit{hash} de integridade à sessão. Em avaliações autorizadas para modalidade remota, quando houver consentimento explícito do usuário, o sistema pode coletar as coordenadas aproximadas via API W3C de Geolocalização, subsidiando investigações de comissões acadêmicas em casos de suspeita fundamentada;
    \item \textbf{Delimitação de Sandbox e Privacidade W3C:} Cabe ressaltar formalmente que, devido às especificações universais de segurança e privacidade dos navegadores modernos mantidas pelo consórcio W3C, aplicações web executadas em modo \textit{sandbox} são estritamente impedidas de acessar identificadores de hardware de camada de enlace (como o endereço físico \textit{MAC Address} da interface de rede). Portanto, a eficácia do modelo de integridade apoia-se deliberadamente na combinação entre barreiras da camada de aplicação (tokens dinâmicos e travas de mutex) e verificações na camada de rede (triagem de IP institucional e canais criptografados TLS).
\end{itemize}

\subsection{Tokens Dinâmicos de Liberação e Janelas Temporais de Aplicação}
\label{sec:seguranca-prova}

O controle temporal e a circunscrição espacial dos exames completam o protocolo de governança da plataforma:
\begin{itemize}
    \item \textbf{Tokens Dinâmicos de Liberação em Sala:} No início da aula presencial de avaliação, a docente projeta ou divulga um código de sessão temporário com validade restrita ao horário letivo, impedindo que discentes iniciem a prova antecipadamente ou a realizem fora do laboratório físico;
    \item \textbf{Temporizador Sincronizado no Servidor:} A contagem regressiva e o encerramento do prazo de submissão baseiam-se estritamente no relógio oficial do servidor, inviabilizando fraudes por manipulação do relógio do sistema operacional do computador do aluno;
    \item \textbf{Janelas Isoladas para Repescagens e Segundas Chamadas:} Para discentes amparados por justificativas acadêmicas, o sistema permite o agendamento de janelas temporais dedicadas vinculadas exclusivamente aos seus respectivos RAs, liberando a avaliação em regime de reposição sem expor o exame à turma regular.
\end{itemize}

\chapter{Implementação e Stack Tecnológica}
\label{cap:implementacao}

Este capítulo trata da engenharia de \textit{software} da plataforma: como ela foi desenvolvida, quais tecnologias foram escolhidas e como as camadas cliente e servidor foram construídas. O objetivo é unir uma simulação local, rápida e fiel, a um servidor concorrente, seguro e de alto desempenho.

Enquanto o Capítulo~\ref{cap:arquitetura} foi organizado segundo o ciclo de vida da prova, este capítulo segue as camadas técnicas do sistema, pois cada componente atende a mais de uma fase. O Quadro~\ref{qua:fases-componentes} indica, para cada fase da prova, onde estão os componentes que a implementam.

\begin{quadro}[htbp]
    \caption{Correspondência entre as fases da prova e os componentes implementados}
    \label{qua:fases-componentes}
    \footnotesize
    \begin{tabularx}{\textwidth}{|>{\raggedright\arraybackslash}p{3.0cm}|L|>{\raggedright\arraybackslash}p{2.6cm}|}
        \hline
        \textbf{Fase da prova} & \textbf{Componentes} & \textbf{Seções} \\
        \hline
        Preparação da turma (Seção~\ref{sec:gestao-academica}) & Modelos GORM de usuários, turmas e matrículas; importação de CSV e envio assíncrono de e-mails & \ref{sec:gorm-persistencia} e \ref{sec:filas-emails-assincronos} \\
        \hline
        Autoria e montagem da prova (Seção~\ref{sec:autoria-conteudo}) & Serializador do VFS para os cenários; MinIO para imagens e \texttt{.tar.gz}; \textit{player} do YouTube para os trechos de vídeo; catálogo de componentes interativos; bloco de teste executado pelo mesmo interpretador do simulador; \textit{seed} do conteúdo atual & \ref{sec:vfs-typescript} e \ref{sec:gorm-persistencia} \\
        \hline
        Abertura da prova (Seção~\ref{sec:abertura-prova}) & Rotas de \textit{token} e janelas no Gin; autenticação JWT; sessão única no Hub WebSocket & \ref{sec:gin-gonic} e \ref{sec:websocket-hub} \\
        \hline
        Resolução (Seção~\ref{sec:validacao-desacoplada}) & Simulador POSIX e VFS no navegador; motor de asserções em \textit{worker pool} no servidor & \ref{sec:vfs-typescript} e \ref{sec:websocket-hub} \\
        \hline
        Supervisão (Seção~\ref{sec:governanca-seguranca}) & Janelas de ajuda e avisos no cliente; Hub WebSocket com eventos para o \textit{Cockpit} & \ref{sec:interface-ts} e \ref{sec:websocket-hub} \\
        \hline
        Encerramento (Seção~\ref{sec:encerramento-prova}) & Registro no IndexedDB e pacote \texttt{.proof}; transações e histórico preservado no PostgreSQL & \ref{sec:visao-stack} e \ref{sec:gorm-persistencia} \\
        \hline
    \end{tabularx}
    \fonte{Autoria própria (2026).}
\end{quadro}

% ----------------------------------------------------------
\section{Metodologia de Engenharia: Adoção Instrumental de Inteligência Artificial e Governança Humana}
\label{sec:metodologia-ia}
% ----------------------------------------------------------

A engenharia de \textit{software} passa por uma mudança estrutural com a chegada de assistentes de IA generativa e de Modelos de Linguagem de Grande Escala (LLMs) aos fluxos de desenvolvimento. Em empresas de tecnologia, o uso desses assistentes como ferramenta já se consolidou para acelerar a produção, prototipar com agilidade e refinar documentação técnica, sem que isso signifique transferir responsabilidade à máquina.

Com base em mais de cinco anos de experiência do autor em liderança técnica no mercado, este trabalho adota essa forma de desenvolvimento acelerado. O uso de IA segue o princípio do \textit{Human-in-the-Loop} (humano no circuito de controle): a ferramenta amplia a capacidade de execução, mas tudo passa por controle humano.

\begin{enumerate}
    \item \textbf{concepção e decisões de domínio:} a ideia do projeto, o modelo de dados, os protocolos criptográficos da contingência \textit{offline}, as regras disciplinares e a dinâmica das avaliações vieram do diagnóstico dos laboratórios da UTFPR e da experiência do autor e da orientadora. Nenhuma regra de negócio, limite de segurança ou decisão de arquitetura foi delegada à máquina;
    \item \textbf{aceleração e apoio à escrita:} o assistente foi usado para automatizar tarefas repetitivas, sugerir rotinas de infraestrutura, transformar conceitos em especificações e ajustar a redação técnica às normas acadêmicas;
    \item \textbf{revisão crítica:} como discutido na Seção~\ref{sec:risco-ia}, o valor do profissional de tecnologia hoje está em auditar o que a ferramenta produz. Cada parágrafo, comando simulado e contrato de API desta monografia foi revisado, testado e validado pelo autor.
\end{enumerate}

Assim, embora a forma do texto e o ritmo de trabalho tenham se beneficiado dessas ferramentas, a concepção, a arquitetura e a responsabilidade científica do trabalho são inteiramente humanas, como já acontece nas práticas mais avançadas da indústria.

% ----------------------------------------------------------
\section{Visão Geral da Stack Tecnológica}
\label{sec:visao-stack}
% ----------------------------------------------------------

As tecnologias foram escolhidas por quatro critérios: escalabilidade, fidelidade da simulação, facilidade de implantação na infraestrutura da universidade e separação clara de responsabilidades. O Quadro~\ref{qua:stack} resume as escolhas.

\begin{quadro}[htbp]
    \caption{Tecnologias adotadas em cada camada}
    \label{qua:stack}
    \footnotesize
    \begin{tabularx}{\textwidth}{|p{2.2cm}|p{3.4cm}|L|}
        \hline
        \textbf{Camada} & \textbf{Tecnologia} & \textbf{Por que foi escolhida} \\
        \hline
        Cliente & TypeScript & Tipagem estática para modelar com precisão as estruturas do POSIX (nós, permissões, credenciais) \\
        \hline
        Cliente & DOM nativo, sem \textit{framework} & Telas escritas como classes TypeScript que montam e desmontam seus próprios elementos, sem dependências externas de interface \\
        \hline
        Cliente & Web Crypto API e IndexedDB & \textit{Hashes} e assinaturas nativos do navegador e armazenamento local transacional de todo o registro da prova, base da contingência \textit{offline} \\
        \hline
        Servidor & Go & Concorrência leve com \textit{goroutines} e canais, ideal para centenas de conexões simultâneas \\
        \hline
        Servidor & Gin-Gonic & Roteamento HTTP rápido, com \textit{middlewares} para autenticação e autorização \\
        \hline
        Servidor & GORM e PostgreSQL & Mapeamento objeto-relacional tipado, transações ACID e exclusão lógica \\
        \hline
        Cliente & YouTube IFrame Player API & Exibição dos trechos de vídeo dos blocos, com início e fim, e leitura da duração para validar o trecho no editor \\
        \hline
        Servidor & MinIO & Armazenamento de objetos para as imagens dos módulos e os arquivos \texttt{.tar.gz} dos cenários \\
        \hline
        Comunicação & WebSockets (WSS) & Canal bidirecional persistente para telemetria e eventos em tempo real \\
        \hline
    \end{tabularx}
    \fonte{Autoria própria (2026).}
\end{quadro}

% ----------------------------------------------------------
\section{Front-end e Núcleo de Simulação (Client-side)}
\label{sec:frontend-simulacao}
% ----------------------------------------------------------

Rodar a simulação no navegador evita tanto as máquinas virtuais quanto um contêiner Docker por aluno, o que seria inviável para laboratórios com orçamento de infraestrutura limitado.

\subsection{Modelagem do VFS e Execução POSIX em TypeScript}
\label{sec:vfs-typescript}

O núcleo da simulação foi escrito inteiramente em TypeScript. A tipagem estática é essencial para representar com precisão as estruturas do POSIX:

\begin{enumerate}
    \item \textbf{nós e metadados (\textit{i-nodes}):} cada entrada do VFS é modelada por interfaces que representam arquivos regulares, diretórios e \textit{links} simbólicos, com permissões em notação octal e simbólica, dono e grupo (UID/GID), referências aos filhos e datas de modificação;
    \item \textbf{permissões e \texttt{umask}:} as permissões \textit{rwx} e a máscara de criação \texttt{umask} são calculadas com operações bit a bit (\textit{bitwise}), reproduzindo o controle de acesso discricionário (DAC) do Linux;
    \item \textbf{troca de usuário:} o simulador mantém uma pilha de sessões que reproduz a troca de credenciais com \texttt{su} e \texttt{sudo}, incluindo as restrições do arquivo \texttt{/etc/sudoers};
    \item \textbf{interpretador de comandos:} o \textit{parser} trata \textit{pipelines} (\texttt{|}), redirecionamentos (\texttt{>}, \texttt{>>}, \texttt{<}) e expansão de variáveis de ambiente, tudo no navegador e sem atraso na digitação.
\end{enumerate}

\subsection{Componentização da Interface em TypeScript}
\label{sec:interface-ts}

A interface é escrita em TypeScript puro, manipulando diretamente o DOM (\textit{Document Object Model}) do navegador, sem \textit{framework} de interface. Cada tela é uma classe que sabe se montar em um contêiner e se desmontar ao sair, e a navegação entre telas é feita pelo endereço da página. Os principais módulos são:

\begin{enumerate}
    \item \textbf{terminal ANSI:} emulador compatível com sequências de escape ANSI e com a paleta do GNOME Terminal, com autocompletar por \texttt{Tab} para comandos e caminhos, histórico pelas setas e os editores \texttt{nano} e \texttt{vim};
    \item \textbf{painel de roteiros e telemetria discente:} painel retrátil com objetivos, enunciados, progresso de questões e cronômetro sincronizado periodicamente com o relógio do servidor;
    \item \textbf{entrada da prova e contingência:} tela para validação do \textit{token} emitido pela professora e captura de pacotes criptográficos em modo desconectado;
    \item \textbf{grade de supervisão do Cockpit docente:} componente reativo em mosaico que renderiza os \textit{cards} discentes atualizados via \texttt{cockpit:telemetry-update}, ordenando dinamicamente a fila de atendimento prioritário (``mãozinha virtual'') e exibindo métricas de tempo de permanência por questão;
    \item \textbf{espelhamento ao vivo (\textit{Live Terminal Mirror}):} painel modal do \textit{Cockpit} que renderiza uma réplica visual síncrona do terminal do aluno e da árvore de arquivos do VFS, alimentado pelo evento \texttt{cockpit:terminal-mirror};
    \item \textbf{janelas de apoio e comunicação direta:} modais para solicitação de ajuda com justificativa obrigatória, telas de advertência escalonada na detecção de \textit{DevTools} e camada sobreposta translúcida para recepção de orientações privadas da professora (\texttt{teacher:direct-message}).
\end{enumerate}

\subsection{Funcionalidades Utilitárias de Perfil}
\label{sec:perfil-utilitario}

A interface também inclui as telas de autoedição de perfil, comuns a alunos e professores: dados de contato, avatar e redefinição de senha com o usuário autenticado. Trata-se de formulários e rotas CRUD convencionais, sem particularidades em relação à prática corrente da indústria, e por isso não são aprofundados aqui. Como explicado na Seção~\ref{sec:modelo-dados}, o detalhamento técnico deste capítulo prioriza o que diferencia o projeto: o núcleo de emulação POSIX e o VFS, a mitigação de fraudes via \textit{DevTools}, a telemetria via \textit{WebSockets} e a resiliência a falhas de rede.

% ----------------------------------------------------------
\section{Back-end e Concorrência de Alta Performance (Server-side)}
\label{sec:backend-golang}
% ----------------------------------------------------------

O servidor é a autoridade confiável da plataforma: guarda as regras de negócio, os dados acadêmicos e garante a integridade das avaliações.

\subsection{A Linguagem Golang e o Modelo de Concorrência CSP}
\label{sec:golang-concorrencia}

A linguagem Go \cite{donovan2015go} foi escolhida pela capacidade de lidar com muita concorrência de rede gastando pouca memória \cite{cox2020golang}:

\begin{enumerate}
    \item \textbf{\textit{goroutines} leves:} \textit{threads} do sistema operacional reservam de 1 a 2\,MB de pilha cada uma, e modelos de \textit{loop} de eventos em uma única \textit{thread} podem travar durante a correção de asserções. Já as \textit{goroutines} começam com pilhas de apenas 2\,KB, que crescem conforme a necessidade. Isso permite manter centenas de conexões \textit{WebSocket} por sala com impacto mínimo na memória do servidor;
    \item \textbf{canais (CSP):} telemetria, fila de dúvidas e despacho de eventos trocam mensagens por canais tipados (\textit{channels}), inspirados no modelo de Processos Sequenciais Comunicantes (\textit{Communicating Sequential Processes} -- CSP) \cite{donovan2015go,cox2020golang}. Isso evita condições de corrida (\textit{race conditions}) e reduz a necessidade de travas manuais (\textit{mutexes}).
\end{enumerate}

\subsection{Roteamento de APIs REST com Gin-Gonic}
\label{sec:gin-gonic}

As APIs seguem o estilo arquitetural REST \cite{fielding2000rest} e são servidas pelo \textit{framework} \textbf{Gin-Gonic}, que usa árvores de prefixo (\textit{radix tree}) para rotear requisições com rapidez e sem alocações desnecessárias de memória:

\begin{enumerate}
    \item \textbf{autenticação e autorização:} \textit{middlewares} verificam os \textit{tokens} JWT, os papéis de usuário (RBAC: Público, Aluno e Professor) e a expiração das sessões;
    \item \textbf{rotas de negócio:} gestão de turmas e matrículas, criação de roteiros de avaliação, emissão e validação dos \textit{tokens} de sala e homologação de notas com justificativa do professor.
\end{enumerate}

\subsection{Persistência e Mapeamento Objeto-Relacional com GORM}
\label{sec:gorm-persistencia}

A persistência de dados utiliza a biblioteca \textbf{GORM} integrada ao PostgreSQL, combinando integridade referencial estrita e modelagem orientada a objetos em Go:

\begin{enumerate}
    \item \textbf{modelos tipados e migrações:} usuários, turmas, matrículas, avaliações, sessões, submissões, observações do professor e eventos da linha do tempo são \textit{structs} Go com \textit{tags} GORM que configuram chaves primárias UUID, índices compostos e integridade referencial;

    \item \textbf{extensão relacional de submissões e sessões:} a entidade de respostas individuais (\texttt{respostas\_questoes\_sessoes}) foi estendida com suporte completo ao fluxo de revisão formativa e auditoria forense. Cada registro possui identificador próprio, chaves estrangeiras para a sessão e a questão (\texttt{sessao\_id} e \texttt{questao\_id}), a pontuação preliminar calculada pelo motor avaliador (\texttt{nota\_sugerida}), a pontuação final confirmada ou revisada pela professora (\texttt{nota\_homologada}), a fundamentação textual para eventuais divergências (\texttt{justificativa\_ajuste}), a devolutiva pedagógica específica (\texttt{feedback\_docente}) e a trilha completa de comandos digitados, tempos de resposta e saídas em coluna JSONB (\texttt{comandos\_executados}, mapeada via \texttt{datatypes.JSON}). Complementarmente, a tabela de sessões de avaliação (\texttt{sessoes\_prova}) incorpora o campo global \texttt{parecer\_geral\_docente}, a nota consolidada (\texttt{nota\_final\_homologada}), o carimbo temporal da homologação (\texttt{homologado\_em}) e o identificador do docente responsável (\texttt{homologado\_por\_id});

    \item \textbf{homologação transacional ACID:} a validação, recálculo e homologação das notas operam estritamente sob transações ACID gerenciadas pelo GORM (\texttt{tx := db.Begin()} \ldots{} \texttt{tx.Commit()}). Dentro do escopo transacional, o serviço atualiza em lote as notas homologadas e os \textit{feedbacks} individuais por questão, recalcula atomicamente a média ou soma ponderada final da avaliação persistindo-a em \texttt{sessoes\_prova}, grava o parecer pedagógico geral e registra um evento irrevogável na linha do tempo de auditoria com o carimbo de tempo da ação. Qualquer inconsistência ou falha de rede aciona reversão integral (\texttt{tx.Rollback()}), impedindo visualizações parciais ou notas descompassadas no portal do estudante;

    \item \textbf{transferência de turma em uma transação:} a troca de turma roda igualmente em uma transação ACID: a matrícula de origem passa a \texttt{TRANSFERIDO} e a nova matrícula é criada ativa na turma de destino em uma única operação. Submissões, notas e comandos antigos continuam ligados ao UUID e ao RA do aluno, sem registros órfãos nem cadastros duplicados;

    \item \textbf{exclusão lógica com preservação de auditoria:} a entidade de usuário usa o tipo \texttt{gorm.DeletedAt}. Nas consultas comuns, o GORM acrescenta automaticamente a condição \texttt{deleted\_at IS NULL}. Para auditorias junto à coordenação da UTFPR, consultas com \texttt{db.Unscoped()} recuperam todo o histórico, inclusive os registros de alunos desligados;

    \item \textbf{suspensão com efeito imediato:} quando o professor marca um aluno como \texttt{INATIVO}, o serviço grava a mudança no banco e envia um sinal, por canal interno, ao Hub de \textit{WebSockets}. As conexões abertas daquele usuário são encerradas na hora, e as próximas requisições à API são barradas pelo \textit{middleware} de autenticação.
\end{enumerate}

\subsection{Hub WebSocket Nativo e Streaming de Telemetria}
\label{sec:websocket-hub}

A comunicação bidirecional em tempo real é sustentada por um \textbf{Hub de WebSockets} nativo desenvolvido em Go, em estrita conformidade com a RFC 6455 \cite{fette2011websocket}:

\begin{enumerate}
    \item \textbf{Hub central e roteamento por sala:} uma \textit{goroutine} central opera como despachante de eventos, gerenciando o registro de novos terminais (\texttt{register}), o encerramento seguro (\texttt{unregister}) e a distribuição de mensagens segmentadas por sala de avaliação;

    \item \textbf{arquitetura de duas \textit{goroutines} por conexão:} cada cliente conectado (discente ou docente) é atendido por um par independente de rotinas assíncronas: uma rotina de leitura contínua (\textit{readPump}) e uma de despacho com amortecimento (\textit{writePump}). Esse desacoplamento previne que eventuais latências de rede ou pausas de renderização no navegador de um estudante comprometam a entrega de pacotes aos demais pares da sala;

    \item \textbf{eventos dedicados de supervisão do Cockpit:} o protocolo define mensagens tipadas em formato JSON para operacionalizar a telemetria em tempo real:
    \begin{itemize}
        \item \texttt{cockpit:telemetry-update}: transporta cargas leves com o total de questões concluídas e pendentes, a identificação da questão corrente e as métricas de temporização (tempo decorrido na tarefa atual e média acumulada por questão);
        \item \texttt{cockpit:terminal-mirror}: transmite em fluxo contínuo os blocos incrementais de caracteres digitados no terminal, códigos de retorno, saídas textuais e o estado serializado da árvore do VFS;
        \item \texttt{teacher:direct-message}: despacha mensagens de intervenção pedagógica direta e orientações privativas (1:1) do docente para a conexão exclusiva do estudante;
    \end{itemize}

    \item \textbf{desacoplamento de goroutines e prevenção de gargalos de I/O no espelhamento:} o espelhamento contínuo da digitação de múltiplos alunos em tempo real possui potencial para saturar o barramento de rede caso implementado via difusão indiscriminada. Para eliminar qualquer risco de contenção de I/O no servidor:
    \begin{itemize}
        \item o Hub adota subscrição seletiva sob demanda (\textit{targeted subscription}): eventos \texttt{cockpit:terminal-mirror} só são despachados aos canais docentes que tenham solicitado ativamente a abertura da tela espelhada daquele discente específico no painel;
        \item as transmissões utilizam canais com \textit{buffer} calibrados para o tráfego em rajadas típico de digitação (\texttt{chan []byte});
        \item emprega-se seleção não bloqueante (\texttt{select} com cláusula \texttt{default}) nas filas de espelhamento transitório. Se o canal de saída de um observador atingir a capacidade máxima por congestionamento temporário de rede, quadros intermediários do espelho são graciosamente descartados sem bloquear a \textit{goroutine} de leitura do aluno nem degradar os canais críticos de notas e batimentos de integridade (\textit{heartbeats});
    \end{itemize}

    \item \textbf{avaliação desacoplada fora do canal de mensageria:} ao término de cada comando prático acompanhado de \textit{snapshot} do sistema de arquivos virtual, a reconstrução e a verificação de asserções são descarregadas para um conjunto de trabalhadores assíncronos (\textit{worker pool}). O canal WebSocket permanece desimpedido para novas interações, assegurando resposta com latência imperceptível para o aluno.
\end{enumerate}

\subsection{Processamento Assíncrono de Filas, Agendamento e Despacho Multicanal}
\label{sec:filas-emails-assincronos}

Importar uma turma por CSV, com dezenas ou centenas de alunos, ou publicar um aviso para a turma gera muitos e-mails. Enviá-los via SMTP dentro da própria requisição HTTP seria inviável: negociar a conexão TLS e enviar cada mensagem leva de 500 milissegundos a 2 segundos por destinatário, o que estouraria o tempo limite (\textit{Gateway Timeout}, HTTP 504) e deixaria a professora esperando.

Por isso, o servidor usa um modelo produtor-consumidor (\textit{producer-consumer}) com um conjunto fixo de trabalhadores (\textit{worker pool}) e um agendador, ilustrado na Figura~\ref{fig:worker-pool}.

\begin{figure}[htbp]
    \centering
    \caption{Envio assíncrono de e-mails com canal e conjunto de trabalhadores}
    \label{fig:worker-pool}
    \begin{tikzpicture}[node distance=0.9cm and 1.1cm]
        \node[box, text width=2.6cm] (csv) {Importação CSV\\(\textit{handler} Gin)};
        \node[box, text width=2.6cm, below=0.5cm of csv] (sched) {Agendador\\(\texttt{time.Ticker})};
        \node[box, text width=2.5cm, right=1.4cm of $(csv)!0.5!(sched)$] (chan) {Canal\\\texttt{chan EmailJob}};
        \node[box, text width=1.9cm, right=1.4cm of chan, yshift=1.0cm] (w1) {\textit{worker} 1};
        \node[box, text width=1.9cm, right=1.4cm of chan] (w2) {\textit{worker} 2};
        \node[box, text width=1.9cm, right=1.4cm of chan, yshift=-1.0cm] (wn) {\textit{worker} $N$};
        \node[box, text width=1.6cm, right=1.1cm of w2] (smtp) {Servidor\\SMTP};

        \draw[arrow] (csv.east) -- (chan);
        \draw[arrow] (sched.east) -- (chan);
        \draw[arrow] (chan) -- (w1.west);
        \draw[arrow] (chan) -- (w2.west);
        \draw[arrow] (chan) -- (wn.west);
        \draw[arrow] (w1.east) -- (smtp);
        \draw[arrow] (w2.east) -- (smtp);
        \draw[arrow] (wn.east) -- (smtp);
        \node[font=\scriptsize, below=0.3cm of sched] {responde \texttt{202 Accepted} na hora};
        \node[font=\scriptsize, below=0.3cm of wn] {limite de taxa e \textit{backoff}};
    \end{tikzpicture}
    \fonte{Autoria própria (2026).}
\end{figure}

O fluxo funciona assim:

\begin{enumerate}
    \item \textbf{fila em canal com \textit{buffer}:} o \textit{handler} HTTP valida os dados (CSV ou novo aviso), grava tudo no banco em uma transação e coloca as tarefas de envio em um canal tipado com \textit{buffer} (\texttt{chan EmailJob}). A resposta \texttt{202 Accepted} volta ao cliente em poucos milissegundos;
    \item \textbf{agendador de avisos:} para avisos programados, uma \textit{goroutine} com temporizador (\texttt{time.Ticker}) consulta periodicamente o PostgreSQL. Quando chega o horário de \texttt{publicar\_em}, o aviso passa a \texttt{PUBLICADO}, o evento \texttt{announcement:new} é enviado aos alunos conectados e os e-mails entram na fila;
    \item \textbf{trabalhadores com limite de taxa:} um número fixo $N$ de \textit{goroutines} consome o canal continuamente. Um limitador de taxa (\textit{rate limiter}) controla o ritmo de envio para respeitar as cotas e as regras antispam do servidor de e-mail da universidade;
    \item \textbf{e-mails em HTML:} cada trabalhador monta a mensagem a partir de modelos HTML com estilos \textit{inline}, compatíveis com os principais leitores de e-mail, com o logotipo da instituição, a identificação da disciplina, o texto do aviso e \textit{links} de ação protegidos por \textit{tokens} de uso único;
    \item \textbf{confirmação de leitura em tempo real:} quando o aluno confirma a leitura de um aviso importante, o servidor grava data, hora e IP e envia \texttt{cockpit:announcement-receipt}, atualizando na hora o percentual de adesão e a lista de pendentes no painel do professor;
    \item \textbf{novas tentativas com espera crescente:} se o servidor SMTP falhar temporariamente (sobrecarga ou erro no TLS), o trabalhador tenta de novo com intervalos cada vez maiores (\textit{exponential backoff}). Esgotadas as tentativas, o erro é registrado em \textit{log} e a notificação fica marcada como pendente, para que a professora possa reenviá-la pelo painel.
\end{enumerate}

\chapter{Avaliação e Resultados}
\label{cap:resultados}

Este capítulo apresenta como a plataforma foi validada e quais resultados foram obtidos até o momento. A Seção~\ref{sec:testes-cobertura} descreve a suíte de testes automatizados que verifica a fidelidade do simulador; a Seção~\ref{sec:resistencia-fraudes} define o ensaio de integridade que avalia a resistência a fraudes; e a Seção~\ref{sec:impacto-pedagogico} compara o custo de preparar o ambiente de aula com máquinas virtuais e com o simulador no navegador.

% ----------------------------------------------------------
\section{Cobertura de Testes e Fidelidade POSIX}
\label{sec:testes-cobertura}
% ----------------------------------------------------------

A fidelidade do simulador é verificada por uma suíte de testes automatizados executada com o \textit{framework} \texttt{vitest}. A maior parte dos casos segue o mesmo roteiro de ponta a ponta: cria-se uma máquina simulada nova, prepara-se o cenário exigido pela questão (usuários, grupos, arquivos e permissões), executa-se a sequência de comandos no interpretador e, ao final, chama-se a função de verificação da questão, que inspeciona o estado do sistema de arquivos virtual. Assim, cada teste confirma ao mesmo tempo que os comandos se comportam como no Linux e que a correção por estado, descrita na Seção~\ref{sec:reconciliacao-servidor}, aceita a resposta.

O Quadro~\ref{qua:suite-testes} descreve os grupos de testes que compõem a suíte.

\begin{quadro}[htbp]
    \caption{Composição da suíte de testes automatizados}
    \label{qua:suite-testes}
    \footnotesize
    \begin{tabularx}{\textwidth}{|p{4.2cm}|L|}
        \hline
        \textbf{Grupo} & \textbf{O que verifica} \\
        \hline
        Desafios oficiais com solução de referência & Conjuntos de desafios práticos (Linux Básico, Linux Médio, Linux Avançado, LPI Linux Essentials, LPIC-1 e Servidor Escola): cada solução de referência é executada no simulador e precisa ser aprovada pela verificação de estado; um caso adicional por conjunto confirma a quantidade de desafios cadastrados \\
        \hline
        Formas alternativas de resolução & Para parte dos desafios, ao menos três soluções diferentes da oficial (caminhos relativos ou absolutos, notação octal ou simbólica, ordens diferentes de comandos) precisam ser aceitas; um caso adicional confirma que todos os cenários apontam para desafios existentes \\
        \hline
        Simulado teórico & Consistência das questões objetivas inspiradas em certificações (alternativas, gabarito e explicação) \\
        \hline
        Conteúdo do tópico de permissões & Presença dos conceitos (\texttt{/etc/passwd}, leitura do \texttt{ls -l}, \textit{rwx}, \texttt{chmod}, \texttt{sudo}) e da lição interativa de concessão de privilégios \\
        \hline
        Cola de comandos & Geração do material de consulta e tratamento de caracteres especiais \\
        \hline
    \end{tabularx}
    \fonte{Autoria própria (2026).}
\end{quadro}

Os grupos de desafios cobrem, na prática, as áreas centrais da disciplina: navegação e manipulação de arquivos e diretórios, permissões \textit{rwx} em notação octal e simbólica, máscara \texttt{umask}, \textit{sticky bit}, usuários e grupos, \texttt{sudo}, redirecionamentos e \textit{pipes}, busca de texto e de arquivos, além das tarefas de administração cobradas nas certificações LPI. O grupo de formas alternativas é o que melhor evidencia o princípio de correção por estado: soluções diferentes da oficial são aceitas porque deixam o sistema de arquivos no estado esperado.

Na versão avaliada, todos os casos de teste foram aprovados. Como a suíte acompanha a evolução do código, sua composição detalhada pode mudar à medida que a plataforma for implementada; o que se mantém é o critério de validação por estado descrito acima.

% ----------------------------------------------------------
\section{Resistência a Fraudes em Ambiente de Prova}
\label{sec:resistencia-fraudes}
% ----------------------------------------------------------

A resistência a fraudes é avaliada por um ensaio de integridade que reproduz, de forma controlada, as tentativas de manipulação que um estudante com conhecimentos de desenvolvimento web poderia fazer durante a prova. O ensaio parte da vulnerabilidade descrita na Seção~\ref{sec:dores-pedagogicas}: no protótipo executado apenas no navegador, as regras de correção ficam acessíveis pelo \textit{DevTools}.

O Quadro~\ref{qua:vetores-ataque} apresenta os vetores de ataque do ensaio e o comportamento esperado em cada arquitetura.

\begin{quadro}[htbp]
    \caption{Vetores de ataque do ensaio de integridade}
    \label{qua:vetores-ataque}
    \footnotesize
    \begin{tabularx}{\textwidth}{|p{3.4cm}|L|L|}
        \hline
        \textbf{Vetor} & \textbf{Protótipo somente no cliente} & \textbf{Arquitetura híbrida} \\
        \hline
        Ler o código de correção no \textit{DevTools} & As condições de aprovação ficam visíveis no código & O cliente recebe apenas \textit{hashes} irreversíveis (Seção~\ref{sec:envelope-cifrado}) \\
        \hline
        Alterar o estado do VFS em memória pelo console & A questão é aprovada sem que os comandos tenham sido executados & O \textit{replay} no servidor não reproduz o estado forjado e a divergência é registrada (Seção~\ref{sec:reconciliacao-servidor}) \\
        \hline
        Forçar a função de verificação a retornar verdadeiro & Aprovação indevida & A nota oficial vem apenas do servidor \\
        \hline
        Editar o registro de comandos guardado no navegador & Sem efeito, pois não há registro oficial & O \textit{checksum} assinado deixa de conferir e o pacote é rejeitado \\
        \hline
        Abrir a prova em um segundo dispositivo & Não detectado & A sessão anterior é encerrada e a colisão aparece no \textit{Cockpit} (Seção~\ref{sec:mutex-sessao}) \\
        \hline
    \end{tabularx}
    \fonte{Autoria própria (2026).}
\end{quadro}

Para cada vetor, o ensaio segue três passos: \textit{(i)} o aluno de teste resolve parte da prova legitimamente; \textit{(ii)} executa a manipulação pelo console do navegador; \textit{(iii)} verifica-se a nota calculada e os registros da linha do tempo de auditoria. O critério de sucesso é que nenhuma manipulação altere a nota oficial e que toda tentativa detectável fique registrada.

No protótipo atual, executado apenas no navegador, os três primeiros vetores são viáveis, já que todo o código de verificação é entregue ao cliente; essa constatação motivou a evolução arquitetural descrita no Capítulo~\ref{cap:arquitetura}. Os resultados do ensaio na arquitetura híbrida serão registrados nesta seção após a conclusão do serviço de retaguarda.

% ----------------------------------------------------------
\section{Impacto Pedagógico e Otimização do Tempo Letivo}
\label{sec:impacto-pedagogico}
% ----------------------------------------------------------

Para estimar o ganho operacional, comparou-se o custo de colocar um aluno diante do primeiro comando em dois cenários: uma máquina virtual tradicional e o simulador no navegador. Os valores do simulador foram medidos; os da máquina virtual são estimativas baseadas na alocação de memória usual e na experiência relatada pela docente da disciplina (Seção~\ref{sec:dores-pedagogicas}).

As medições do simulador foram feitas no Google Chrome 152, em modo \textit{headless}, sobre o \textit{build} de produção servido localmente, em um computador com 4 núcleos e 8\,GB de RAM. Foram realizadas cinco execuções com o \textit{cache} desativado; em cada uma, a página de exercícios foi aberta e 20 comandos foram digitados no terminal. A memória total da aba foi obtida pelo PSS (\textit{proportional set size}) do processo renderizador no sistema operacional. A Tabela~\ref{tab:comparativo-recursos} resume os resultados.

\begin{table}[htbp]
    \centering
    \caption{Recursos necessários para iniciar a prática: máquina virtual e simulador}
    \label{tab:comparativo-recursos}
    \footnotesize
    \begin{tabular}{p{5.6cm}cc}
        \toprule
        \textbf{Indicador} & \textbf{Máquina virtual} & \textbf{Simulador no navegador} \\
        \midrule
        Memória dedicada ao ambiente & 2\,048 a 4\,096\,MB\textsuperscript{a} & 105\,MB (aba inteira)\textsuperscript{b} \\
        Memória além de uma aba vazia & --- & 68\,MB\textsuperscript{b} \\
        Memória usada pelo código do simulador (\textit{heap} JavaScript) & --- & 2,2\,MB\textsuperscript{b} \\
        Download ou imagem necessária & Imagem ISO de alguns GB\textsuperscript{a} & 1,4\,MB\textsuperscript{b} \\
        Tempo até o primeiro comando & Dezenas de minutos\textsuperscript{a} & 0,9 a 1,9\,s\textsuperscript{b} \\
        Instalação e privilégios administrativos & Necessários & Não necessários \\
        \bottomrule
    \end{tabular}
    \fonte{Autoria própria (2026). \textsuperscript{a}Valores estimados. \textsuperscript{b}Valores medidos.}
\end{table}

Mesmo considerando a aba inteira, incluindo o motor do navegador, o simulador consome cerca de 20 a 40 vezes menos memória do que a faixa usualmente reservada a uma máquina virtual, e o código do simulador em si ocupa cerca de 2\,MB. O tempo até o primeiro comando cai de dezenas de minutos para cerca de um segundo, o que é coerente com a perda de 30 a 45 minutos por aula relatada na Seção~\ref{sec:dores-pedagogicas}. A medição do tempo de início em aulas reais, com turmas da disciplina, fica como etapa seguinte da avaliação.

\chapter{Conclusão e Trabalhos Futuros}
\label{cap:conclusao}

Este capítulo resume as contribuições técnicas e pedagógicas do trabalho, discute suas limitações e aponta caminhos para continuá-lo.

\section{Considerações Finais}
\label{sec:consideracoes-finais}
A plataforma proposta cumpre o objetivo de democratizar o acesso à prática de Sistemas Operacionais: elimina as barreiras da virtualização e oferece ao professor um ambiente robusto, seguro e com métricas precisas de aprendizagem.

\section{Trabalhos Futuros}
\label{sec:trabalhos-futuros}
Como continuidade deste trabalho, propõe-se:
\begin{enumerate}
    \item ampliar o conjunto de comandos POSIX e incluir utilitários de rede, como \texttt{tcpdump}, \texttt{curl} e \texttt{iptables};
    \item integrar a plataforma ao Moodle da UTFPR por meio do protocolo LTI (\textit{Learning Tools Interoperability});
    \item usar técnicas de inteligência artificial para identificar padrões de dificuldade e recomendar desafios de fixação adaptados a cada aluno.
\end{enumerate}


% ----------------------------------------------------------
% ELEMENTOS PÓS-TEXTUAIS
% ----------------------------------------------------------
\postextual

\bibliography{pos-textual/referencias}

\end{document}
